% GENERATED FILE. Edit canonical lessons, then run npm run generate:books.
% canonical-source-hash: fnv1a64:1370579991fbaf1f
% canonical-lessons: TE-C491-arigipovu, TE-C491-bujupattu, TE-C491-velisipovu, TE-C491-pagilipovu, TE-C491-ksinincu

\chapter{Wear and Tear}
\label{ch:te-wear-and-tear}
\hlchaptermodality{491}

\begin{chapteropening}
\textbf{By the end of this chapter:} \emph{I can say to wear down, to go mouldy, to fade, to crack and to deteriorate.}

Complete the last lesson of chapter 491: I can say to wear down, to go mouldy, to fade, to crack and to deteriorate.
\end{chapteropening}

\section[\texorpdfstring{\te{అరిగిపోవు} (arigipōvu)}{అరిగిపోవు (arigipōvu)}]{\te{అరిగిపోవు} (arigipōvu) — to wear down; to wear out}

\label{lesson:TE-C491-arigipovu}



\begin{quote}
Before the new one: say the Telugu for to split, then the Telugu for to break off.
\end{quote}

\subsection*{You'll want to know: \te{అరిగిపోవు}}
\textbf{\te{అరిగిపోవు}} — \emph{arigipōvu} — \textquotedblleft{}to wear down; to wear out\textquotedblright{}.

\begin{quote}
\te{ఎక్కువ} \te{నడవడంతో} \te{నా} \te{చెప్పులు} \te{అరిగిపోయాయి}.
\end{quote}

\begin{quote}
\emph{ekkuva naḍavaḍantō nā ceppulu arigipōyāyi.}
\end{quote}

\begin{quote}
With so much walking, my sandals have worn down.
\end{quote}

Past (it …): \te{అరిగిపోయింది}, \emph{arigipōyindi}.

Said of soles, tyres and steps worn thin by use. Its subject is a thing, so the past shown is the third person: \te{అరిగిపోయింది}, it wore down.

\begin{cousinweb}
\textbf{\te{అరిగిపోవు}} (\emph{arigipōvu}) is native Dravidian: DEDR 316 lists Telugu \emph{aṟugu}, \textquotedblleft{}to wear away, to waste away; to be digested\textquotedblright{}, here with \te{పోవు}, \textquotedblleft{}to go\textquotedblright{} (DEDR 4572), which marks a change that is complete.
\end{cousinweb}

\begin{grammarlens}[title={What you've built}]
One more word for getting around.
\end{grammarlens}

\subsection*{Your turn}
\begin{itemize}
\raggedright
  \item \emph{Say it:} \emph{arigipōvu}
  \item \emph{Say it:} \emph{arigipōvu}, once more
  \item \emph{Say it:} \emph{ekkuva naḍavaḍantō nā ceppulu arigipōyāyi.} — With so much walking, my sandals have worn down.
  \item \emph{Recall:} say the Telugu for to split, then the Telugu for to break off, then say \emph{arigipōvu} again
\end{itemize}

\begin{tcolorbox}[breakable,colback=teal!4,colframe=teal!35!black,title={Before you move on}]
What does \te{అరిగిపోవు} mean? (\textquotedblleft{}To wear down; to wear out\textquotedblright{}.) Say it once more.
\end{tcolorbox}
\section[\texorpdfstring{\te{బూజుపట్టు} (būjupaṭṭu)}{బూజుపట్టు (būjupaṭṭu)}]{\te{బూజుపట్టు} (būjupaṭṭu) — to go mouldy}

\label{lesson:TE-C491-bujupattu}



\begin{quote}
Before the new one: say the Telugu for to break off, then the Telugu for to wear down.
\end{quote}

\subsection*{You'll want to know: \te{బూజుపట్టు}}
\textbf{\te{బూజుపట్టు}} — \emph{būjupaṭṭu} — \textquotedblleft{}to go mouldy\textquotedblright{}.

\begin{quote}
\te{తేమ} \te{వల్ల} \te{రొట్టె} \te{బూజుపట్టింది}.
\end{quote}

\begin{quote}
\emph{tēma valla roṭṭe būjupaṭṭindi.}
\end{quote}

\begin{quote}
Because of the damp the bread went mouldy.
\end{quote}

Past (it …): \te{బూజుపట్టింది}, \emph{būjupaṭṭindi}.

\te{బూజు}, mould, + \te{పట్టు}, to catch. A house cobweb is also \te{బూజు}. Its subject is a thing, so the past shown is the third person: \te{బూజుపట్టింది}, it went mouldy.

\begin{cousinweb}
\textbf{\te{బూజుపట్టు}} (\emph{būjupaṭṭu}) is native Dravidian: \te{బూజు}, \textquotedblleft{}mould, mildew\textquotedblright{} (DEDR 4357), and \te{పట్టు}, \textquotedblleft{}to catch\textquotedblright{} (DEDR 4034).
\end{cousinweb}

\begin{grammarlens}[title={What you've built}]
One more word for getting around.
\end{grammarlens}

\subsection*{Your turn}
\begin{itemize}
\raggedright
  \item \emph{Say it:} \emph{būjupaṭṭu}
  \item \emph{Say it:} \emph{būjupaṭṭu}, once more
  \item \emph{Say it:} \emph{tēma valla roṭṭe būjupaṭṭindi.} — Because of the damp the bread went mouldy.
  \item \emph{Recall:} say the Telugu for to break off, then the Telugu for to wear down, then say \emph{būjupaṭṭu} again
\end{itemize}

\begin{tcolorbox}[breakable,colback=teal!4,colframe=teal!35!black,title={Before you move on}]
What does \te{బూజుపట్టు} mean? (\textquotedblleft{}To go mouldy\textquotedblright{}.) Say it once more.
\end{tcolorbox}
\section[\texorpdfstring{\te{వెలిసిపోవు} (velisipōvu)}{వెలిసిపోవు (velisipōvu)}]{\te{వెలిసిపోవు} (velisipōvu) — to fade}

\label{lesson:TE-C491-velisipovu}



\begin{quote}
Before the new one: say the Telugu for to wear down, then the Telugu for to go mouldy.
\end{quote}

\subsection*{You'll want to know: \te{వెలిసిపోవు}}
\textbf{\te{వెలిసిపోవు}} — \emph{velisipōvu} — \textquotedblleft{}to fade\textquotedblright{}.

\begin{quote}
\te{ఎండకు} \te{చీర} \te{రంగు} \te{వెలిసిపోయింది}.
\end{quote}

\begin{quote}
\emph{eṇḍaku cīra raṅgu velisipōyindi.}
\end{quote}

\begin{quote}
In the sun the colour of the sari faded.
\end{quote}

Past (it …): \te{వెలిసిపోయింది}, \emph{velisipōyindi}.

Said of colours and of old photos and memories. Its subject is a thing, so the past shown is the third person: \te{వెలిసిపోయింది}, it faded.

\begin{cousinweb}
\textbf{\te{వెలిసిపోవు}} (\emph{velisipōvu}) is a Telugu compound whose first part has an origin that is not settled; the second is the native \te{పోవు}, \textquotedblleft{}to go\textquotedblright{} (DEDR 4572).
\end{cousinweb}

\begin{grammarlens}[title={What you've built}]
One more word for getting around.
\end{grammarlens}

\subsection*{Your turn}
\begin{itemize}
\raggedright
  \item \emph{Say it:} \emph{velisipōvu}
  \item \emph{Say it:} \emph{velisipōvu}, once more
  \item \emph{Say it:} \emph{eṇḍaku cīra raṅgu velisipōyindi.} — In the sun the colour of the sari faded.
  \item \emph{Recall:} say the Telugu for to wear down, then the Telugu for to go mouldy, then say \emph{velisipōvu} again
\end{itemize}

\begin{tcolorbox}[breakable,colback=teal!4,colframe=teal!35!black,title={Before you move on}]
What does \te{వెలిసిపోవు} mean? (\textquotedblleft{}To fade\textquotedblright{}.) Say it once more.
\end{tcolorbox}
\section[\texorpdfstring{\te{పగిలిపోవు} (pagilipōvu)}{పగిలిపోవు (pagilipōvu)}]{\te{పగిలిపోవు} (pagilipōvu) — to crack; to burst}

\label{lesson:TE-C491-pagilipovu}



\begin{quote}
Before the new one: say the Telugu for to go mouldy, then the Telugu for to fade.
\end{quote}

\subsection*{You'll want to know: \te{పగిలిపోవు}}
\textbf{\te{పగిలిపోవు}} — \emph{pagilipōvu} — \textquotedblleft{}to crack; to burst\textquotedblright{}.

\begin{quote}
\te{చేతి} \te{నుంచి} \te{పడి} \te{గ్లాసు} \te{పగిలిపోయింది}.
\end{quote}

\begin{quote}
\emph{cēti nuñci paḍi glāsu pagilipōyindi.}
\end{quote}

\begin{quote}
The glass fell from my hand and broke.
\end{quote}

Past (it …): \te{పగిలిపోయింది}, \emph{pagilipōyindi}.

Said of glass, pots, pipes and lips in the cold. Its subject is a thing, so the past shown is the third person: \te{పగిలిపోయింది}, it broke.

\begin{cousinweb}
\textbf{\te{పగిలిపోవు}} (\emph{pagilipōvu}) is native Dravidian: DEDR 3808 lists Telugu \emph{pagulu}, \textquotedblleft{}to break, to crack, to burst\textquotedblright{}, here with \te{పోవు}, \textquotedblleft{}to go\textquotedblright{} (DEDR 4572).
\end{cousinweb}

\begin{grammarlens}[title={What you've built}]
One more word for getting around.
\end{grammarlens}

\subsection*{Your turn}
\begin{itemize}
\raggedright
  \item \emph{Say it:} \emph{pagilipōvu}
  \item \emph{Say it:} \emph{pagilipōvu}, once more
  \item \emph{Say it:} \emph{cēti nuñci paḍi glāsu pagilipōyindi.} — The glass fell from my hand and broke.
  \item \emph{Recall:} say the Telugu for to go mouldy, then the Telugu for to fade, then say \emph{pagilipōvu} again
\end{itemize}

\begin{tcolorbox}[breakable,colback=teal!4,colframe=teal!35!black,title={Before you move on}]
What does \te{పగిలిపోవు} mean? (\textquotedblleft{}To crack; to burst\textquotedblright{}.) Say it once more.
\end{tcolorbox}
\section[\texorpdfstring{\te{క్షీణించు} (kṣīṇiñcu)}{క్షీణించు (kṣīṇiñcu)}]{\te{క్షీణించు} (kṣīṇiñcu) — to deteriorate; to decline}

\label{lesson:TE-C491-ksinincu}



\begin{quote}
Before the new one: say the Telugu for to fade, then the Telugu for to crack.
\end{quote}

\subsection*{You'll want to know: \te{క్షీణించు}}
\textbf{\te{క్షీణించు}} — \emph{kṣīṇiñcu} — \textquotedblleft{}to deteriorate; to decline\textquotedblright{}.

\begin{quote}
\te{వృద్ధుడి} \te{ఆరోగ్యం} \te{క్షీణించింది}.
\end{quote}

\begin{quote}
\emph{vṛddhuḍi ārōgyaṁ kṣīṇiñcindi.}
\end{quote}

\begin{quote}
The old man's health deteriorated.
\end{quote}

Past (it …): \te{క్షీణించింది}, \emph{kṣīṇiñcindi}.

Said of health, of an old building and of a river that runs low. Its subject is a thing, so the past shown is the third person: \te{క్షీణించింది}, it declined.

\begin{cousinweb}
\textbf{\te{క్షీణించు}} (\emph{kṣīṇiñcu}) is built from Sanskrit \emph{kṣīṇa}, \textquotedblleft{}diminished, wasted away\textquotedblright{} (Monier-Williams), with the Telugu verb ending -\te{ఇంచు}.
\end{cousinweb}

\begin{grammarlens}[title={What you've built}]
One more word for getting around.
\end{grammarlens}

\subsection*{Your turn}
\begin{itemize}
\raggedright
  \item \emph{Say it:} \emph{kṣīṇiñcu}
  \item \emph{Say it:} \emph{kṣīṇiñcu}, once more
  \item \emph{Say it:} \emph{vṛddhuḍi ārōgyaṁ kṣīṇiñcindi.} — The old man's health deteriorated.
  \item \emph{Recall:} say the Telugu for to fade, then the Telugu for to crack, then say \emph{kṣīṇiñcu} again
\end{itemize}

\begin{tcolorbox}[breakable,colback=teal!4,colframe=teal!35!black,title={Before you move on}]
What does \te{క్షీణించు} mean? (\textquotedblleft{}To deteriorate; to decline\textquotedblright{}.) Say it once more.
\end{tcolorbox}
